% part: intuitionistic-logic
% chapter: propositions-as-types
% section: introduction

\documentclass[../../../include/open-logic-section]{subfiles}

\begin{document}

\olfileid{int}{pty}{int}

\olsection{引言}

λ演算是一个简单的计算模型。它操作由变元构成的项。如果 $N$ 和 $M$ 是项，那么 $(NM)$ 也是项（“$N$ 应用于~$M$”）。如果 $N$ 是项，那么 $\lambd[x][N]$ 也是项。如果 $N$ 包含变元~$x$，那么在 $\lambd[x][N]$ 中相应的 $x$ 的出现就被 $\lambd[x]$ 算子所约束，因此不再是自由的。形如 $(\lambd[x][N]M)$ 的项被称为可 $\beta$-归约为 $\Subst{N}{M}{x}$（即将 $N$ 中每一个自由出现的 $x$ 替换为 $M$ 所得的结果）。如果 $N'$ 是由 $N$ 对某个子项进行 $\beta$-归约得到的，则称项 $N$ 可 $\beta$-转换为 $N'$，记作 $N \redone N'$。λ演算中的计算是一个序列 $N \redone N_1
\redone N_2 \redone \dots$。如果该序列终止于项 $N_k$，且其中没有任何子项可以再进行 $\beta$-归约，则称该项为\emph{正规的}（normal），或称其为~$N$ 的一个范式。不妨将λ演算中的计算视为这样的序列。如果计算得到一个范式，则计算停机，该范式即为计算的结果。结果表明，这个简单的计算模型是图灵完备的：每一个部分递归函数都可以由λ演算中的一个项来计算。

Haskell Curry（柯里）和 William Howard（霍华德）各自独立地发现了自然演绎与λ演算之间的密切相似性。直观上看，项 $\lambd[x][N]$ 是一个函数：$x$ 是自变元；$N$ 告诉我们在给定~$x$ 时如何计算其值。那么项 $(\lambd[x][N]M)$ 就是将函数 $\lambd[x][N]$ 应用于自变元~$M$，而 $\beta$-归约告诉我们如何对其进行求值：用 $M$ 替换 $N$ 中的 $x$（并继续对所得的项求值）。现在，将它们视为具有固定定义域和值域的函数。如果 $\lambd[x][N]$ 的定义域是 $A$，值域是~$B$，那么我们应当认为变元~$x$ 只接受来自~$A$ 的自变元，而 $N$ 则产生属于~$B$ 的值。因此，仅当 $\lambd[x][N]$ 是 $A \to B$ 的函数且 $M \in A$ 时，$(\lambd[x][N]M)$ 才是有意义的。如果我们想把这种信息加入λ演算，就得到了\emph{类型化}λ演算。在类型化λ演算中，并非每个表达式 $(\lambd[x][N]M)$ 都是合法的，只有那些 $\lambd[x][N]$ 和 $M$ 的“类型”相匹配的表达式才是合法的，即 $\lambd[x][N]$ 具有类型 $A \lif B$ 且 $M$ 具有类型~$A$。$\lambd[x][N]$ 的类型为 $A \to B$，仅当 $x$ 的类型为 $A$ 且 $N$ 的类型为~$B$。一般而言，并非每个表达式 $(M_1M_2)$ 都合法。只有那些 $M_1$ 具有类型 $A \to B$ 且 $M_2$ 具有类型~$A$ 的项 $(M_1M_2)$ 才合法；若如此，则 $(M_1M_2)$ 具有类型~$B$。

这些类型规则清晰地对应于自然演绎中的推导，其中类型由公式表示，而推导本身则对应于λ项。例如，$!A \lif !B$ 的一个推导对应于项 $\lambd[\typeof{x}{!A}][N]$，该项具有函数类型 $!A \lif !B$。自然演绎的推理规则对应于类型化λ演算中的类型规则，例如，

\begin{prooftree}
  \AxiomC{$\Discharge{!A}{x}$}
  \DeduceC{$!B$}
  \DischargeRule{\Intro\lif}{x}
  \UnaryInfC{$!A \lif !B$}
  \DisplayProof\bottomAlignProof
  \quad\text{对应}\quad
  \Axiom$x: !A \fCenter N: !B$
  \UnaryInf$ \fCenter \lambd[\typeof{x}{!A}][N] : !A \lif !B$
\end{prooftree}

其中右边的规则表示：如果 $x$ 的类型是~$!A$ 且 $N$ 的类型是~$!B$，那么 $\lambd[\typeof{x}{!A}][N]$ 的类型就是~$!A \lif !B$。

$\Elim\lif$ 规则对应于复合项的类型规则，即
\[
  \AxiomC{$!A \lif !B$}
  \AxiomC{$!A$}
  \RightLabel{\Elim\lif}
  \BinaryInfC{$!B$}
  \DisplayProof
  \quad\text{对应}\quad
  \Axiom$\fCenter M_1: !A \lif !B$
  \Axiom$\fCenter M_2: !A$
  \BinaryInf$ \fCenter (M_1M_2) : !B$
  \DisplayProof
\]
如果一个 \Intro\lif{} 规则之后紧跟着一个 \Elim\lif{} 规则，该推导就可以被简化：
\begin{gather*}
  \AxiomC{$\Discharge{!A}{x}$}
  \DeduceC{$!B$}
  \DischargeRule{\Intro\lif}{x}
  \UnaryInfC{$!A \lif !B$}
  \AxiomC{}
  \DeduceC{$!A$}
  \RightLabel{\Elim\lif}
  \BinaryInfC{$!B$}
  \DisplayProof
  \quad\redone\quad
  \AxiomC{}
  \DeduceC{$!A$}
  \DeduceC{$!B$}
  \DisplayProof
\end{gather*}
这对应于λ项的 $\beta$-归约：
\[
(\lambd[\typeof{x}{!A}][N]M) \quad\redone\quad
\Subst{N}{M}{x}.
\]

类似的对应关系也存在于 $\land$ 的规则与“积”类型之间，以及 $\lor$ 的规则与“和”类型之间。

简单类型化λ演算中的项与自然演绎推导之间的这种对应关系，被称为“Curry--Howard”对应、“证明即程序”或“命题即类型”对应。除了公式（命题）对应于类型、证明对应于项之外，我们可以将对应关系总结如下：

\begin{center}
  \begin{tabular}{c c c}
    逻辑 & 程序 \\
    \hline
    命题/公式 & 类型 \\
    证明 & 项 \\
    假设 & 变元 \\
    被解除的假设 & 约束变元 \\
    未被解除的假设 & 自由变元 \\
    $!A \lif !B$ & 函数类型 \\
    $!A \land !B$ & 积类型 \\
    $!A \lor !B$ & 和类型 \\
    $\lfalse$ & 空类型 \\
    \hline
  \end{tabular}
\end{center}

Curry--Howard 对应是自动证明助理和程序类型检查器的基石之一，因为检查一个见证命题的证明（正如我们上面所做的那样），就相当于检查一个程序（项）是否具有所声明的类型。
\end{document}